\begin{onehalfspace}

Questo capitolo riassume gli strumenti software e l'ambiente usati nella parte sperimentale. La motivazione estesa di ciascuna scelta --- il confronto tra i framework, i metodi di simulazione di Aer, le specifiche GPU e la configurazione WSL --- è riportata nell'Appendice~\ref{app:strumenti}.

\textbf{Framework: Qiskit + Qiskit Aer.} L'intera pipeline usa \textbf{Qiskit} \cite{qiskit2023}, l'SDK di IBM Quantum, e il suo simulatore \textbf{Qiskit Aer}. La scelta è motivata dalla coerenza con l'hardware superconduttore IBM su cui è calibrato il modello di rumore, dalla maturità del noise model di Aer rispetto alle alternative (Cirq, PennyLane), e dall'allineamento con la letteratura e con l'implementazione di riferimento dell'algoritmo di Shor \cite{skosana2021, github_shor}.

\textbf{Metodo di simulazione e limite di scala.} Tra i metodi di \texttt{AerSimulator} si adotta \texttt{statevector} (simulazione esatta del vettore di stato), per la precisione numerica richiesta dagli use case a 11--18 qubit. Il costo in memoria cresce come $2^n$ (ampiezze \texttt{complex128}): con 12\,GB di memoria il limite pratico è $\approx 30$ qubit. È questo vincolo esponenziale a rendere impraticabile la simulazione diretta dello Shor logico con surface code e a imporre l'uso di Stim per i circuiti \emph{stabilizer} della QEC (Capitoli~\ref{chap:surface} e~\ref{chap:integrazione}).

\textbf{Strumenti della parte QEC.} La stima del \textit{logical error rate} del surface code richiede milioni di esecuzioni di circuiti stabilizer, fuori portata per un simulatore statevector. Si adottano quindi due strumenti specializzati: \textbf{Stim} \cite{stim2021} (versione 1.16.0), che campiona \textit{detection events} di circuiti Clifford sfruttando il teorema di Gottesman--Knill, e \textbf{PyMatching} \cite{pymatching2021} (versione 2.4.0), che implementa la decodifica \ac{MWPM} a partire dal modello di errore del circuito. La divisione dei compiti --- Qiskit Aer per i circuiti generali di piccola taglia, Stim e PyMatching per la \ac{QEC} stabilizer --- è discussa nel Capitolo~\ref{chap:integrazione}.

\textbf{Modello di rumore.} Il modulo \texttt{qiskit\_aer.noise} compone errore \textbf{depolarizzante}, \textbf{rilassamento termico} $T_1/T_2$ e \textbf{readout error} (matrice di confusione), con parametri calibrati sull'hardware \texttt{ibm\_marrakesh} \cite{epjqt_noise2024}; la configurazione dettagliata è nel Capitolo~\ref{chap:obiettivi}. Prima dell'esecuzione i circuiti sono transpilati (\texttt{optimization\_level=2}) per minimizzare profondità e accumulo di errori.

\textbf{Machine learning e ambiente.} Il classificatore del Metodo 2 è addestrato con \textbf{scikit-learn} \cite{scikit_learn}, scelto per l'interfaccia unificata \texttt{fit}/\texttt{predict} che consente di confrontare più modelli sullo stesso dataset. Le simulazioni girano in \textbf{WSL2} (Ubuntu 22.04) su \ac{CPU}: l'accelerazione \ac{GPU} tramite \texttt{qiskit-aer-gpu} (NVIDIA RTX 4070) era prevista dal piano iniziale ma non è risultata utilizzabile in questo ambiente --- il motivo tecnico e le conseguenze sui tempi di esecuzione sono documentati nella Sezione~\ref{sec:gpu_limit}. Su questa base, il Capitolo~\ref{chap:metodologia} descrive l'architettura del pipeline sperimentale.

\end{onehalfspace}
